The channel protocol records authorization, each write attempt, its returned byte count,
complete receiver requests, and response receipt separately (details in Appendix~\ref{app:socket}).

The forced ordering test produced confirmed post-acknowledgement byte acceptance in 20/20 early-check attempts, versus 0/20 fenced attempts. The authorization marker preceded acknowledgement in these early-check races, exposing handoffs that the stream-wrapper authorization-marker estimand would miss. This forced schedule does not estimate its natural frequency.

There were 80 crash-point attempts. In the 20 native-return-before-log cases, 20 retained unknown handoff counts; absence of a count was never converted to zero. Receiver prefix logs, restart ordering, collection-epoch retry and re-grant outcomes are preserved separately.

Under receiver backpressure (800 attempts; Table~\ref{tab:revision-latency} in Appendix~\ref{app:socket}), the median withdrawal acknowledgement with per-chunk fencing stayed at or below 0.8\,ms, while a lock held over the whole write reached 4.9\,s with slow or stalled receivers; 0 withdrawals went unacknowledged.

\paragraph{Matched metadata.} Five variants share durable lineage, epochs, copied job extras
and the same HTTP transport. Each of seven schedules has 20 valid trials per AVD configuration (140 per configuration), with all five variants paired in each trial. Four read-only instances per configuration run disjoint seeded trial ranges on the same host; these are not independent devices. Table~\ref{tab:matched}
reports confirmed violations and ambiguous requests separately; neither unknowns nor delivery deficits
are counted as successful suppression. Tenants within each trial remain paired in uncertainty analysis.
\begin{table}[t]\centering\caption{Matched-metadata results. TT: affected tenant-trials; V: confirmed violating TT; U: requests with unknown handoff/content evidence. Counts aggregate two AVD configurations; uncertainty resamples complete paired trials within configuration, schedule and instance. All instances use the host renderer with Vulkan disabled.}
\label{tab:matched}\mtab
\begin{tabular}{@{}lrrr@{}}\toprule\mhead Variant & TT & V & U\\\midrule

enqueue & 1400 & \mbad{1038} & 3 \\

job start & 1400 & \mbad{236} & 3 \\

layered early & 1400 & \mgood{0} & 6 \\

chunk fence & 1400 & \mgood{0} & 2 \\

layered + fence & 1400 & \mgood{0} & 1 \\

\bottomrule\end{tabular}\end{table}

The analysis excludes 8 failed paired trials and retains their diagnostics. Among the fixed permitted-flow controls (T3, T4 backup, T6, T7), layered early leaves 168/14000 permitted items undelivered, chunk fence leaves 96/14000 permitted items undelivered, layered + fence leaves 60/14000 permitted items undelivered. All 448 such items across the five variants come from reboot trials in which JobScheduler stopped an upload at shutdown and the fixture declined rescheduling; no affected trace shows a policy block or send failure (Appendix~\ref{app:matched-trials}).

The matched layered-early versus fence contrast and trial-level intervals are reported below.

AVD A: paired confirmed-violation difference (layered early minus fence) 0.00 percentage points, percentile interval [0.00, 0.00], across 140 complete paired trials.

AVD B: paired confirmed-violation difference (layered early minus fence) 0.00 percentage points, percentile interval [0.00, 0.00], across 140 complete paired trials.

Layered early checks also have no confirmed violations in these schedules; the enlarged campaign does not demonstrate a natural-schedule advantage for final placement. A zero-event percentile interval is descriptive, not a population zero-risk bound. The forced HTTP and TLS barriers isolate the narrower residual race. Unknowns remain separate from confirmed violations.

At the trial level, layered early, chunk fence and layered + fence have no confirmed violation in 280 paired trials each; counting unknowns as violations, the one-sided 95\% bounds stay at or below 4.2\% if trials were independent (Appendix~\ref{app:matched-trials}).
